\documentclass{article}
\usepackage[utf8]{inputenc}
\usepackage[T1]{fontenc}
\usepackage{geometry}
\geometry{
    a4paper,
    left=2.5cm,
    right=2.5cm,
    top=2.5cm,
    bottom=2.5cm
}
\usepackage{graphicx}
\usepackage{float}
\usepackage{hyperref}
\usepackage[english]{babel}
\usepackage{amsmath}
\usepackage{amssymb}
\usepackage{array}
\usepackage{booktabs}
\usepackage{caption}
\usepackage{subcaption}
\usepackage{enumitem}
\usepackage{fancyhdr}
\usepackage{titlesec}
\usepackage{url}
\usepackage{textcomp}
\usepackage{xcolor}
\usepackage{listings}
\usepackage{multirow}
\usepackage{longtable}
\usepackage{tabularx}
\usepackage{times}

% Set up page style
\pagestyle{fancy}
\fancyhf{}
\rfoot{\thepage}
\renewcommand{\headrulewidth}{0pt}

% Title formatting
\titleformat{\chapter}[display]{\normalfont\bfseries}{\chaptertitle\thechapter}{0pt}{\Huge}
\titlespacing*{\chapter}{0pt}{50pt}{40pt}

% Document info
\title{Template Document}
\author{Academic Support Agent}
\date{\today}

\begin{document}
\maketitle

%Define some colors
\definecolor{w}{rgb}{1,1,1}
\definecolor{b}{rgb}{0,0,0}

%Title Slide
{\usebackgroundtemplate{\includegraphics[width=\paperwidth]{Fondo1.png}} %Title Background
\begin{frame}
\rmfamily %Title Font
\color{w}\bfseries\titlepage %Title Color & Weight
\end{frame}
}

%Contents Slide
{\usebackgroundtemplate{\includegraphics[width=\paperwidth]{Fondo2.png}} %Contents Background
\begin{frame}
\frametitle{ ndice} %Contents Title
\rmfamily %Contents Font
\begin{itemize} %Start enumeration
\color{w} %Define Text Color  
\item 1. Introducci n
\item 2. Primer apartado
\item 3. Segundo, etc
\end{itemize}
\end{frame}
}

%Normal Slide (copy, paste and modify this slide for longer presentations)
\begin{frame}
\frametitle{1. Introducci n} %Title
\framesubtitle{Una descripci n m s detallada, si se quiere} %Subtitle
\rmfamily %Font
\color{b} %Color
Aqu  se puede incluir texto, aunque quiz  es m s pr ctico dividirlo en puntos:
\begin{itemize}
\color{b}
\item Por ejemplo, as ,
\item O as . Podemos incluir una f rmula
\end{itemize}

\begin{equation*}
f(P,C) := \sum_{\boldsymbol{p}\in{P}}\displaystyle{\minimize_{\boldsymbol{c}\in C} \parallel\boldsymbol{p} - \boldsymbol{c}\parallel^2}
\end{equation*}

Y con la funci n \textit{includegraphics} un gr fico o imagen.
\end{frame}
\end{document}